\documentclass[11pt]{letter}
\usepackage[margin=1in]{geometry}
\usepackage[colorlinks=true,urlcolor=blue]{hyperref}

\signature{Muhammad R.~Hasyim\\
Simons Center for Computational Physical Chemistry\\
New York University\\
\href{mailto:mh7373@nyu.edu}{mh7373@nyu.edu}}

\address{Muhammad R.~Hasyim\\
Simons Center for Computational Physical Chemistry\\
New York University\\
New York, NY 10003}

\begin{document}

\begin{letter}{Editors\\
\textit{Nature Computational Science}}

\opening{Dear Editors,}

Please consider the enclosed Brief Communication, ``Path Sampling Diagnoses
Protein Diffusion Models,'' for publication in \textit{Nature Computational
Science}.

This manuscript introduces a path-sampling diagnostic for explicit-noise
generative structure models. Because the random variables used during reverse
diffusion can be recorded and assigned exact prior probabilities, the full
denoising trajectory can be sampled with transition path sampling. The
manuscript then combines this trajectory ensemble with on-the-fly probability
enhanced sampling to expose rare ligand-pose basins that ordinary single-seed
inference misses.

The Brief Communication format fits the scope: one rare-event sampling
framework, one high-profile generative model and one protein--ligand case
study. The biological demonstration uses Boltz-2 on the MEK1--FZC complex, a
system that connects directly to current debates about memorization and
generalization in protein--ligand co-folding models. The TPS+OPES ensemble
contains a dominant native-like basin together with persistent secondary
ligand-pose basins above 1~\AA\ RMSD, showing that model uncertainty and
alternative denoising outcomes can be characterized at the path level.

The manuscript should interest readers working at the boundary of statistical
mechanics, enhanced sampling, generative modeling and computational structural
biology.
The main text is intentionally compact, while the Supplementary Information
contains the path-probability derivations, quotient-space diffusion context,
evaluation metrics and reproducible computational workflow needed by specialist
readers.

The work is not under consideration elsewhere. Related manuscript and software
materials are available in the accompanying
\href{https://github.com/muhammadhasyim/tps-diffusion-model}{project repository}.
I have had no prior discussions with a \textit{Nature Computational Science}
editor about this work. Large language model assistance was used for manuscript
restructuring and prose drafting; all scientific claims, analysis choices and
final text were reviewed by the author.

\closing{Sincerely,}

\end{letter}

\end{document}
